%!TeX program = xelatex
\documentclass[zihao=-4]{ctexart}  % 小四号字

\usepackage[a4paper, left=2.5cm, right=2.5cm, top=2.5cm, bottom=2.5cm]{geometry}
\usepackage{tcolorbox}
\tcbuselibrary{skins}
\usepackage{enumitem}
\usepackage{fontspec}

% 字体设置（Windows 用 SimSun，macOS 请改为 Songti SC）
\setCJKmainfont{SimSun}[AutoFakeBold=2.5]
\setmainfont{Times New Roman}

\pagestyle{empty}          % 封面不显示页码
\setlist{nosep, leftmargin=*, labelsep=0.5em}   % 紧凑列表

\begin{document}

% ========== 表格外部（居中，带横线填空） ==========
\begin{center}
    % 标题：黑体小二加粗
    {\heiti\zihao{-2}\textbf{华东理工大学 XXXX---XXXX 学年第 X 学期}}

    \vspace{0.5\baselineskip}

    % “课程论文”行
     {\heiti\zihao{-2}\textbf\textbf{《XX》课程论文 \quad XXXX.X}}

    % 课程论文行下方行距
    \vspace{1.5\baselineskip}

    % 班级行
    班级 \underline{\makebox[3cm][c]{XXXX}} \quad 学号 \underline{\makebox[3cm][c]{XXXXXXXX}} \quad 姓名  \underline{\makebox[4cm][c]{XXXX}}

    % 班级行下方行距
    \vspace{1\baselineskip}

    % 开课学院行
    开课学院  \underline{\makebox[3cm][c]{XXXX}} \quad 任课教师 \underline{\makebox[4cm][c]{XXXX}} \quad 成绩 \underline{\hspace{2cm}}
\end{center}

%\vspace{1.2\baselineskip}

% ========== 表格内部（三个无间隙子边框） ==========
% 所有框内标题设为宋体四号，枚举项设为宋体五号

% ----- 框1：论文题目（一行，无空白）-----
\begin{tcolorbox}[
    sharp corners,
    boxrule=0pt,
    leftrule=0.8pt, rightrule=0.8pt,
    toprule=0.8pt, bottomrule=0.8pt,
    left=4pt, right=4pt, top=4pt, bottom=4pt,
    before skip=0pt, after skip=0pt,
    colback=white
]
    % 论文题目标题：宋体四号（不加粗）
    \noindent{\zihao{4} 论文题目：TEST-TEST} \hspace{4cm}
    \vspace{0pt}
\end{tcolorbox}

% ----- 框2：论文要求 -----
\begin{tcolorbox}[
    sharp corners,
    boxrule=0pt,
    leftrule=0.8pt, rightrule=0.8pt,
    toprule=0pt, bottomrule=0.8pt,
    left=4pt, right=4pt, top=4pt, bottom=4pt,
    before skip=0pt, after skip=0pt,
    colback=white
]
    % 论文要求标题：宋体四号
    \noindent{\zihao{4} 论文要求：}
    % 枚举列表：宋体五号
    \begin{enumerate}[label=\arabic*., leftmargin=2em]
        \item {\zihao{5} line1；}
        \item {\zihao{5} line2；}
        \item {\zihao{5} line3；}
        \item {\zihao{5} line4；}
        \item {\zihao{5} line-end。}
    \end{enumerate}
\end{tcolorbox}

% ----- 框3：教师评语（无下边框）-----
\begin{tcolorbox}[
    sharp corners,
    boxrule=0pt,
    leftrule=0.8pt, rightrule=0.8pt,
    toprule=0pt, bottomrule=0pt,
    left=4pt, right=4pt, top=4pt, bottom=0pt,
    before skip=0pt, after skip=0pt,
    colback=white
]
    % 教师评语标题：宋体四号
    \noindent{\zihao{4} 教师评语：}
    \vspace{4cm}   % 留出手写评语的空间
    \hfill
\end{tcolorbox}

% ----- 框4：签字日期（带下划线签字区）-----
\begin{tcolorbox}[
    sharp corners,
    boxrule=0pt,
    leftrule=0.8pt, rightrule=0.8pt,
    toprule=0pt, bottomrule=0.8pt,
    left=4pt, right=4pt, top=4pt, bottom=4pt,
    before skip=0pt, after skip=0pt,
    colback=white
]
    \hfill   % 将内容推至右侧
    \raggedleft
    % 签字日期字体使用五号（也可保持与全局一致）
    {\zihao{5} 教师签字： \underline{\hspace{3cm}} \\
    \vspace{0.8\baselineskip}
    年 \hspace{1cm} 月 \hspace{1cm} 日}
\end{tcolorbox}

\end{document}